\mask{53}{358}{237}{201}
\mask{53}{562}{237}{80}
\mask{323}{416}{236}{228}
\mask{55.7883}{83.4}{36.8419}{12.52}
\mask{94.0302}{83.4}{50.1961}{12.52}
\mask{145.626}{83.4}{410.685}{12.52}
\mask{145.626}{97.32}{410.685}{12.28}
\mask{145.626}{111}{410.685}{23.32}
\mask{145.626}{135.72}{410.685}{12.04}
\mask{145.626}{149.16}{410.685}{11.92}
\mask{145.626}{162.48}{410.685}{34.96}
\mask{145.626}{198.84}{410.685}{12.04}
\mask{145.626}{212.28}{410.685}{23.56}
\mask{145.626}{237.24}{410.685}{11.92}
\mask{145.626}{250.56}{410.685}{12.07}
\mask{145.626}{264.03}{410.685}{12.04}
\mask{145.626}{277.47}{410.685}{23.44}
\mask{145.626}{302.31}{410.685}{23.56}
\mask{145.626}{327.27}{410.685}{11.92}
\zhbox{54}{64}{504}{17}{14}{\textbf{引脚功能（续）}}
\attext{57}{85}{34}{9}{\centering\textbf{引脚}}
\attext{99}{85}{41}{9}{\centering\textbf{名称}}
\attext{148}{85}{406}{9}{\textbf{说明}}
\attext{148}{99}{406}{9}{\textbf{升压高侧 MOSFET（HS-FET）栅极驱动。}连接 SWD 栅极。}
\attext{148}{112.5}{406}{9}{\textbf{升压高侧 MOSFET 驱动自举电源。}在 BST2 与 SW2 之间接电容。VCC 或 BST1 可向 BST2 供电。}
\attext{148}{137.5}{406}{9}{\textbf{升压低侧 MOSFET（LS-FET）栅极驱动。}连接 SWC 栅极。}
\attext{148}{151}{406}{9}{\textbf{功率地。}栅极驱动电流回流引脚。}
\attext{148}{164.5}{406}{9}{\textbf{驱动电路及内部偏置电源。}使用不小于 2.2 $\mu$F 的陶瓷电容对 VCC 去耦，并尽可能靠近 VCC。VOUT 或 VIN 可向 VCC 供电。}
\attext{148}{201}{406}{9}{\textbf{降压低侧 MOSFET 栅极驱动引脚。}连接 SWB 栅极。}
\attext{148}{214.5}{406}{9}{\textbf{降压高侧 MOSFET 驱动自举电源。}在 BST1 与 SW1 之间接电容。VCC 或 BST2 可向 BST1 供电。}
\attext{148}{239}{406}{9}{\textbf{降压高侧 MOSFET 栅极驱动。}直接连接 SWA 栅极。}
\attext{148}{252.5}{406}{9}{\textbf{变换器降压开关节点。}连接 SWA 源极与 SWB 漏极。}
\attext{148}{266}{406}{9}{\textbf{电源及电压检测输入。}}
\attext{148}{279.5}{406}{9}{\textbf{使能。}不用时接输入电源，以自动启动；也可配置输入欠压锁定（UVLO）保护。EN 不得悬空。}
\attext{148}{304.5}{406}{9}{\textbf{输入电压调节。}设置开关工作期间的最小输入电压。不用时接 AVDD。}
\attext{148}{328.5}{406}{9}{\textbf{I\textsuperscript{2}C 从地址设置。}ADDR 还设置 ENPWR 位的默认值。}
\zhbox{54}{360}{234}{198}{9.5}{\setlength{\parskip}{1bp}\textbf{绝对最大额定值$^{(1)}$}\par\medskip
VIN、EN：$-0.3$ V 至 $+40$ V\par
VOUT、IAVGP、IAVGN：$-0.3$ V 至 $+30$ V\par
VCC：$-0.3$ V 至 $+8.5$ V\par
SW1、SW2：$-1$ V 至 $+40$ V\par
\hspace{1em}（小于 20 ns 时：$-5$ V 至 $+45$ V）\par
LG1、LG2：$-0.3$ V 至 $+10$ V\par
\hspace{1em}（小于 20 ns 时：$-2$ V 至 $+11$ V）\par
BST1、HG1：$-0.3$ V 至 $V_{SW1}+8.5$ V\par
BST2、HG2：$-0.3$ V 至 $V_{SW2}+8.5$ V\par
其他引脚：$-0.3$ V 至 $+6.5$ V\par
连续功耗$^{(2)(5)}$：4.4 W\par
结温：150 $^\circ$C\par
引线温度：260 $^\circ$C\par
存储温度：$-65$ $^\circ$C 至 $+150$ $^\circ$C}
\zhbox{54}{563}{234}{78}{9.5}{\textbf{推荐工作条件$^{(3)}$}\par\medskip
启动电压（$V_{SU}$）：6 V 至 36 V\par
工作输入电压（$V_{IN}$）$^{(4)}$：5 V 至 36 V\par
输出电压（$V_{OUT}$）：0.5 V 至 28 V\par
工作结温（$T_J$）：$-40$ $^\circ$C 至 $+125$ $^\circ$C}
\zhbox{324}{361}{130}{13}{10}{\textbf{热阻}}
\zhbox{324}{417}{234}{226}{7.5}{\textbf{注：}\par
1）超过这些额定值可能损坏器件。\par
2）最大允许功耗取决于最大结温 $T_J(\mathrm{MAX})$、结到环境热阻 $\theta_{JA}$ 和环境温度 $T_A$。任意环境温度下的最大允许连续功耗为 $P_D(\mathrm{MAX})=[T_J(\mathrm{MAX})-T_A]/\theta_{JA}$。超过最大允许功耗会使芯片温度过高，器件可能进入热关断。内部热关断电路可防止器件永久损坏。\par
3）不保证器件在其工作条件之外正常运行。\par
4）该值为输出已稳压至不低于 5 V，且 VCC 负载小于 10 mA 后的工作输入电压。\par
5）在 MPS 评估板 EV2984-R-00B 上测得；板尺寸 91.4 mm $\times$ 66 mm，六层 PCB，铜厚依次为 2 oz、1 oz、1 oz、1 oz、1 oz、2 oz。\par
6）表中 $\theta_{JA}$ 仅适用于与其他封装比较，不得用于设计。这些值按 JESD51-7 在指定 JEDEC 板上计算、仿真，不代表实际应用中的性能。}
